% ════════════════════════════════════════════════════════════════════
\section{nền}

\begin{frame}{lý thuyết trò chơi hợp tác — bảy khái niệm, một ví dụ}
  ba đoạn $a, b, c$. $v(S)$ là điểm khi cho model đọc đúng tập $S$.

  \vspace{0.4em}
  \begin{columns}[t]
  \begin{column}{0.52\textwidth}
    \begin{dia}
v(rỗng)      0,05      ngưỡng tau_q = 0,84
v({a})       0,12
v({b})       0,18
v({c})       0,71      <- cao mà KHÔNG chứa dữ kiện
v({a,b})     0,94      <- bắc cầu
v({a,c})     0,74
v({b,c})     0,76
v({a,b,c})   0,95
    \end{dia}
  \end{column}
  \begin{column}{0.44\textwidth}
    \footnotesize
    \begin{tabular}{@{}p{0.34\linewidth}p{0.6\linewidth}@{}}
      \hl{trò chơi} & $(N, v)$, $N$ là pool, $v$ là hàm giá trị\\
      \hl{đơn điệu} & $S \subseteq S' \Rightarrow v(S) \le v(S')$. \no{ta KHÔNG có}\\
      \hl{thắng} & $v(S) \ge \tau_q$\\
      \hl{mwc} & thắng, và không tập con thật sự nào thắng\\
      \hl{pivotal} & thêm $i$ biến thiếu thành đủ\\
      \hl{banzhaf} & đếm số lần pivotal\\
      \hl{pgi} & tỉ lệ mwc có chứa $i$\\
    \end{tabular}
  \end{column}
  \end{columns}

  \vspace{0.5em}
  \begin{keyb}
    $v$ \no{không đơn điệu}: thêm đoạn gây nhiễu làm $v$ \emph{giảm}.
    định nghĩa mwc không cần đơn điệu, và bao đơn điệu
    $\bar v(S) = \max_{S'\subseteq S} v(S')$ cho $\mathrm{mwc}(\bar v) = W$ —
    đơn điệu hóa \no{không đổi gì ngoài chi phí.}
  \end{keyb}
\end{frame}

\begin{frame}{integrated gradients — và vì sao tầng b cần nó}
  gradient tại một điểm chỉ nói độ dốc \emph{ở đó}. ig tích phân dọc đường từ mốc tới điểm thật:

  \[
    \mathrm{ig}_i \;=\; (x_i - x_i^{0}) \times
    \frac{1}{M}\sum_{m=1}^{M}
    \frac{\partial F\!\left(x^{0} + \alpha_m (x - x^{0})\right)}{\partial x_i},
    \qquad \alpha_m = \tfrac{m}{M}
  \]

  \vspace{0.3em}
  \begin{columns}[t]
  \begin{column}{0.48\textwidth}
    \begin{bx}[title={tính chất ta dùng}]
      \textbf{completeness}\;
      $\sum_i \mathrm{ig}_i = F(x) - F(x^0)$
      \vspace{0.3em}

      tổng đóng góp \emph{bằng đúng} phần mà đầu vào tạo ra. nên credit chia được, và
      kiểm được bằng một đẳng thức.
    \end{bx}
  \end{column}
  \begin{column}{0.48\textwidth}
    \begin{bx}[title={vì sao ig cho cạnh, ablation cho đoạn}]
      \[ \text{ig rẻ hơn ablation} \iff N > 3M \]
      \vspace{-0.6em}
      \begin{dia}
cạnh đồ thị   hàng nghìn, đồ thị CỐ ĐỊNH   -> ig
token         hàng trăm, prompt CỐ ĐỊNH    -> ig
đoạn trong    prompt PHÌNH theo số đoạn    -> ablation
pool          -> tỉ lệ = 3M, K triệt tiêu,
                 KHÔNG có điểm hòa vốn
      \end{dia}
    \end{bx}
  \end{column}
  \end{columns}
\end{frame}

% ════════════════════════════════════════════════════════════════════
\section{tầng a}

\begin{frame}{đối tượng cần tìm}
  \[
    \boxed{\;
      S \text{ đủ} \;\iff\; \underbrace{v(S) \ge \tau_q}_{\text{trả lời đúng}}
      \;\wedge\; \underbrace{\mathrm{adq}(S)}_{\text{có bằng chứng}}
    \;}
  \]
  \[
    S \text{ tối tiểu} \;\iff\; S \text{ đủ} \;\wedge\;
    \nexists\, S' \subsetneq S \text{ đủ}
    \qquad\qquad
    W = \{\, S \text{ đủ tối tiểu} \,\}
  \]

  \vspace{0.6em}
  \begin{columns}[t]
  \begin{column}{0.48\textwidth}
    \begin{bx}[title={dư thừa hiện ra thành bội số mwc}]
      ba đoạn trùng ý $\Rightarrow$ ba mwc khác nhau, mỗi tập chứa một đoạn.
      \vspace{0.3em}

      nên bản này \no{không cần gom cụm để hạ} $2^K$ — dư thừa tự lộ trong cấu trúc $W$,
      còn $2^K$ thì surrogate lo.
    \end{bx}
  \end{column}
  \begin{column}{0.48\textwidth}
    \begin{bx}[title={cái giá so với shapley}]
      mất \textbf{completeness}: $\sum_c \varphi_c = F(E) - F(E')$ là đẳng thức
      chính xác, bắt được mọi lỗi cài đặt.
      \vspace{0.3em}

      mwc không có đẳng thức đó. đổi lại: được cả \emph{họ} lời giải thích, chỉ số
      cấu trúc, và độ tin cậy.
    \end{bx}
  \end{column}
  \end{columns}
\end{frame}

\begin{frame}{«đủ» có hai thành phần, và $v$ chỉ đo một}
  \begin{columns}[t]
  \begin{column}{0.5\textwidth}
    \begin{badb}[title={ca mồi tham số}]
      \begin{dia}
q  ai đạo diễn phim thắng oscar 1994?
y  Robert Zemeckis

đoạn 19  "Forrest Gump là hiện tượng
          văn hóa thập niên 1990"
         -> KHÔNG nói đạo diễn
         -> KHÔNG nói thắng giải

model thấy tên phim -> TỰ NHỚ
                    -> trả lời ĐÚNG
v({19}) = 0,71  cao
      \end{dia}
      \vspace{0.2em}
      định nghĩa chỉ-dùng-$v$ \no{nhận đoạn 19 vào $T$}, rồi dạy retriever đi lấy nó.
    \end{badb}
  \end{column}
  \begin{column}{0.46\textwidth}
    \begin{bx}[title={và lệch này đã được đo}]
      \footnotesize
      \begin{tabular}{@{}p{0.52\linewidth}r@{}}
        model trả lời \emph{đúng} khi ngữ cảnh \emph{không đủ} & 35--62\,\%\\[0.3em]
        tỉ lệ trả lời sai: không ngữ cảnh & 10,2\,\%\\
        \qquad ngữ cảnh không đủ & 66,1\,\%\\[0.3em]
        ca \emph{có} ngữ cảnh đủ sau truy hồi &\\
        \quad musique-ans & 46,2\,\%\\
        \quad hotpotqa & 44,6\,\%\\
      \end{tabular}
      \vspace{0.4em}

      \dd{joren et al., iclr 2025 — arXiv:2411.06037}
      \vspace{0.3em}

      hơn \emph{nửa} truy vấn trên đúng hai bộ ta nhắm \no{không có tập đủ trong pool.}
    \end{bx}
  \end{column}
  \end{columns}

\end{frame}

\begin{frame}{hai điều kiện, bốn ô — mỗi ô một nguyên nhân khác}
  \vspace{0.5em}
  \begin{center}
  \begin{tabular}{@{}>{\raggedright\arraybackslash}p{0.12\linewidth}
                     >{\raggedright\arraybackslash}p{0.37\linewidth}
                     >{\raggedright\arraybackslash}p{0.37\linewidth}@{}}
    \toprule
    & \textbf{adq phủ hết} & \textbf{adq không phủ}\\
    \midrule
    \hl{$v \ge \tau_q$}
      & \yes{nhãn tốt} — nhận
      & \no{mồi tham số} — loại\\[0.7em]
    \hl{$v < \tau_q$}
      & bằng chứng \emph{có} mà model không nối được\newline
        $\rightarrow$ lỗi \textbf{model sinh}, tầng b \no{không sửa được}
      & pool \emph{thiếu}\newline
        $\rightarrow$ thoát sớm ở \texttt{a0}/\texttt{c0}\\
    \bottomrule
  \end{tabular}
  \end{center}

  \vspace{0.9em}
  \begin{keyb}
    giữ $\mathrm{adq}$ \hl{tách khỏi} $v$, không gộp thành một điểm số — gộp thì mất bảng
    này, và mất luôn khả năng tách \emph{lỗi truy hồi} khỏi \emph{lỗi model sinh}.
    ô dưới-trái là ô mà bản chỉ-dùng-$v$ không bao giờ thấy.
  \end{keyb}
\end{frame}

\begin{frame}{$v(S)$ — hai lượt llm, một con số}
  \begin{columns}[t]
  \begin{column}{0.52\textwidth}
    \begin{bx}[title={bước 1 · sinh từ đúng $S$}]
      \begin{dia}
Answer using ONLY the passages below.
If insufficient, say: INSUFFICIENT.

[1] <nội dung chunk thứ nhất trong S>
[2] <nội dung chunk thứ hai>
Question: <q>
Answer:
      \end{dia}
      \vspace{-0.2em}
      \[ \hat a_S = \mathrm{generate}(q, S) \]
    \end{bx}
    \begin{bx}[title={bước 2 · judge, đọc logprob token đầu}]
      \begin{dia}
Question: <q>
Gold answer: Robert Zemeckis
Also acceptable: Zemeckis · Robert Lee ...
Candidate answer: <â_S>
Correct? One word: Yes or No.
Answer:
      \end{dia}
    \end{bx}
  \end{column}
  \begin{column}{0.44\textwidth}
    \[
      v(S) \;=\;
      \frac{\sum_{t \in \mathcal{Y}} e^{\ell_t}}
           {\sum_{t \in \mathcal{Y}} e^{\ell_t} + \sum_{t \in \mathcal{N}} e^{\ell_t}}
      \;\in [0,1]
    \]
    \dd{$\ell_t$ = logprob token đầu · $\mathcal{Y}$ = \{Yes, yes, ...\}}

    \vspace{0.5em}
    \begin{bx}[title={ba chỗ không được làm sai}]
      \footnotesize
      \no{không} log-likelihood của $y$ — \emph{"Paris"} vs \emph{"Paris, France"} đúng
      ngữ nghĩa mà logprob thấp, và lệch đó \emph{không} triệt tiêu giữa các tập con:
      \textbf{thứ tự lật}.
      \vspace{0.35em}

      \no{không} thang 1--5 — 5 mức làm hàng chục tập con hòa nhau, không phân biệt
      được $S$ với $S \setminus \{c\}$, và \emph{tối tiểu} mất nghĩa.
      \vspace{0.35em}

      \hl{phải} nhét \texttt{answer\_aliases} vào prompt judge, không thì tập \emph{đúng}
      bị loại oan.
    \end{bx}
    \vspace{-0.3em}
    \dd{$\mathrm{cost}(S) = \text{overhead} + \sum_{k \in S} \mathrm{len}(k)$ —
        không phải $500 + 250|S|$, vì document ${\sim}100$ và summary ${\sim}600$
        là hai lớp độ dài}
  \end{column}
  \end{columns}
\end{frame}

\begin{frame}{$\mathrm{adq}$ — chẻ câu hỏi, rồi đối chiếu}
  \begin{columns}[t]
  \begin{column}{0.46\textwidth}
    \begin{bx}[title={bước 1 · chẻ thành mảnh}]
      \footnotesize
      \hl{hỏi ai:} musique gán sẵn phân rã 1-hop $+$ đáp án trung gian.
      corpus riêng thì llm, \textbf{1 lượt mỗi truy vấn} — không phải mỗi tập con.
      \begin{dia}
mảnh 1  phim thắng oscar 1994 là Forrest Gump
mảnh 2  Forrest Gump do Zemeckis đạo diễn
      \end{dia}
    \end{bx}
    \begin{bx}[title={bước 2 · đối chiếu}]
      \footnotesize
      \hl{hỏi ai:} model nli nhỏ (${\sim}400$M), chạy local, \textbf{0 lượt api}.
      \[ e(k,j) = \mathrm{nli}(\underbrace{\text{đoạn } k}_{\text{premise}},
                               \underbrace{\text{mảnh } j}_{\text{hypothesis}}) \]
      \dd{$40 \times J$ lượt forward model nhỏ}
    \end{bx}
  \end{column}
  \begin{column}{0.5\textwidth}
    \begin{bx}[title={bảng phủ — 0 lượt llm}]
      \begin{dia}
            mảnh 1   mảnh 2
đoạn  3     entail      —
đoạn  7     entail      —
đoạn 12       —      entail
đoạn 23       —      entail
đoạn 31       —      entail
đoạn 19       —         —      <- bẫy
      \end{dia}
      \vspace{-0.3em}
      \begin{align*}
        \mathrm{cov}(S) &= \{\, j : \exists k \in S,\; e(k,j) = \text{entail} \,\}\\
        \mathrm{adq}(S) &\iff \mathrm{cov}(S) = \{1, \dots, J\}
      \end{align*}
    \end{bx}
    \begin{keyb}
      \footnotesize
      bảng phủ cho \textbf{ba} thứ, cả ba miễn phí:
      \begin{dia}
1 CỬA THOÁT  mảnh nào trống -> bỏ truy vấn ngay
2 ỨNG VIÊN   phủ tập nhỏ, giải CHÍNH XÁC
             {3,12} {3,23} {7,12} ...
3 PRIOR |T|  |T| nên bằng J
      \end{dia}
    \end{keyb}
  \end{column}
  \end{columns}
\end{frame}

\begin{frame}{vì sao phải chẻ ra — hai lý do độc lập}
  \begin{columns}[t]
  \begin{column}{0.48\textwidth}
    \begin{bx}[title={độ dài}]
      premise $=$ cả $S$ có thể tới $10{,}000$ token, mà model nli chỉ nhận $512$.
      \vspace{0.4em}

      chẻ ra thì mỗi lượt kiểm là \textbf{1 đoạn} (${\sim}100$ token) đối
      \textbf{1 câu ngắn}.
    \end{bx}
  \end{column}
  \begin{column}{0.48\textwidth}
    \begin{badb}[title={bắc cầu}]
      claim ghép \emph{"phim thắng 1994 do Zemeckis đạo diễn"} thì
      \no{không đoạn nào một mình} khẳng định được.
      \vspace{0.4em}

      không chẻ $\Rightarrow$ loại oan cả $\{3,12\}$, tức loại đúng
      \textbf{ca multi-hop} mà cả dự án nhắm.
    \end{badb}
  \end{column}
  \end{columns}

  \vspace{0.7em}
  \begin{keyb}
    chẻ ra thì mỗi mảnh có \emph{đúng một} đoạn khẳng định nó, và ghép lại thành $T$ đầy đủ.
    \dd{\;· cơ chế theo arXiv:2608.00585 — kiểm từng đoạn đối câu hỏi gốc thất bại ở multi-hop}
  \end{keyb}
\end{frame}

\begin{frame}{ngưỡng và trần}
  \vspace{-0.2em}
  \begin{align*}
    \tau_q &= \theta \,\bigl[\, v_{\mathrm{ref}}(q) - v(\emptyset) \,\bigr] + v(\emptyset),
      \qquad \theta = 0{,}95\\[-0.1em]
    v_{\mathrm{ref}}(q) &= \max\Bigl(\, v(C),\; v(R),\; \max_k v(\{k\}),\;
      \text{mọi } v(S) \text{ đã quan sát} \Bigr)
  \end{align*}
  \vspace{-0.6em}
  \begin{columns}[t]
  \begin{column}{0.48\textwidth}
    \begin{bx}[title={mốc phải là cực đại quan sát được}]
      khi $C$ còn đoạn nhiễu thì một tập con có thể ghi điểm \emph{cao hơn} cả $C$ —
      neo vào $v(C)$ là neo vào \no{mốc bị dìm}.
      \begin{dia}
v(C)  = 0,81   nhồi 40 đoạn -> phân tán
v(R)  = 0,88   12 đoạn, ít nhiễu hơn
->  v_ref = 0,88,  tau_q = 0,84
      \end{dia}
    \end{bx}
  \end{column}
  \begin{column}{0.48\textwidth}
    \begin{bx}[title={hai mốc đặc biệt}]
      \begin{dia}
v(rỗng)  không cho đoạn nào
         -> chỉ tri thức tham số của model
v(R)     12 đoạn: top-2 mỗi cụm, VÀ tôn
         trọng hạn ngạch nguồn 20/15/5
      \end{dia}
      \vspace{0.2em}
      \dd{\textbf{chốt chặn mẫu số nhỏ:}\;
          $v_{\mathrm{ref}} - v(\emptyset) < 0{,}05 \Rightarrow T = \emptyset$,
          cờ \emph{"truy hồi không đóng góp"} — ca có thật, không phải lỗi.}
    \end{bx}
  \end{column}
  \end{columns}
\end{frame}

\begin{frame}{cập nhật $v_{\mathrm{ref}}$ tự nhất quán — và $\theta < 1$ gánh tất}
  quan sát mới có thể \emph{cao hơn} trần đang dùng. lúc đó $v_{\mathrm{ref}}$ tăng,
  nên $\tau_q$ cũng tăng — liệu tập vừa nâng trần có bị chính ngưỡng mới loại?

  \vspace{0.8em}
  \[
    v(S) > v_{\mathrm{ref}}
    \;\Longrightarrow\;
    v_{\mathrm{ref}} := v(S)
    \;\Longrightarrow\;
    \tau_q' = \theta\bigl[v(S) - v(\emptyset)\bigr] + v(\emptyset)
    \;\underbrace{<}_{\theta\,<\,1}\; v(S)
  \]

  \vspace{0.8em}
  \begin{columns}[t]
  \begin{column}{0.48\textwidth}
    \begin{keyb}[title={không bao giờ tự phủ định}]
      tập vừa nâng trần \yes{luôn vẫn vượt ngưỡng mới}. cập nhật không thể loại
      chính cái gây ra nó.
    \end{keyb}
  \end{column}
  \begin{column}{0.48\textwidth}
    \begin{keyb}[title={không phải xét lại lịch sử}]
      $\tau_q$ chỉ \emph{tăng}, nên mọi tập đã bị loại vẫn bị loại.
      \vspace{0.25em}

      \dd{$\theta = 1$ thì suy biến: $\tau_q' = v(S)$, tập nâng trần chỉ \emph{bằng}
          ngưỡng, và tính chất trên mất}
    \end{keyb}
  \end{column}
  \end{columns}
\end{frame}
